\chapter{Conclusion}

\section{Summary}

\subsection{Methods}

We evaluated how directly applicable published (trace) element ratios are for the determination of the
depositional environment of sediments~\parencite{wei-2020,Worden_1996}. We used XRF, \ICPMS and IRMS to
obtain concentrations of the elements of interest. The uncalibrated XRF method proved too 
unreliable to obtain usable results.
The determination of Cl/Br proved difficult with \ICPMS. The published method for measuring halogens
\parencite{he_determination_2019} could not be reproduced successfully and even led to significant corrosion of the instrument. 
In contrast, the measurement of elements like Sr, Ba, and S worked reliably with the established
HF digestion method, and IRMS produced reliable results for \TOC.

We found that the thresholds for Sr/Ba published by \textcite{wei-2020} do not match the classification based
on the geological map (\textcite{Linth_Map}, \cref{fig:study_area}), while the S/TOC measurements work
with the published thresholds, especially those by \textcite{berner-raiswell-1983}.
Calibrating the thresholds to local conditions improves the results for Sr/Ba.
This calibration
requires certified standards of both depositional environments. For our study such certified
samples were not available, consequently we relied on the classification based on the geological map, 
excluding the obvious outlier JWL~1.
Because we relied on the existing classification to judge if the measured ratios give a reliable classification,
we cannot in turn use the obtained classifications to check if the existing classification is correct.

Overall, the S/TOC ratio provides a strong signal with our data. Sr/Ba separates the samples far less clearly,
and only with calibration. Therefore, it should be considered only when standards for calibration are
available, and as an auxiliary signal.

\subsection{Geology}

We were interested in whether some of the units mapped as \USM in our study area (\cref{fig:study_area})
contain marine layers or may be of marine origin altogether. 
Our data does not allow us to answer that question conclusively, but it strongly indicates that the sample JWL~1
originates from a marine layer.

JWL~1 was collected from a layer with clear evidence of tectonic deformation, and the entire subalpine molasse has been 
deformed and uplifted in the study area. This opens the possibility of the layer having been moved to its current
location by tectonics, rather than having been deposited in-situ. Our study focussed only on the
geochemical analysis of the samples, and cannot make statements about the geology of the study area.


\section{Outlook}

\subsection{Methods}

Throughout the study, the methods we used raised several questions for further inquiry.

\XRF did not yield reliable quantitative results, because the instrument was not calibrated and no standards to
calibrate individual measurements were available.
The instrument should be calibrated before any further measurements are conducted. 
Certified, powdered standards for the elements we are interested in (B, Br, Cl, Ga, S, Sr, Ba) are not readily available.
Therefore, a custom standard would have to be made.

The Cl and Br concentrations measure with \ambif\ digestion confirm that halogens are retained in the digestion cake
and transferred into solution with ammonia~\cref{fig:halogen-sensitivity}, otherwise the
measured counts should be close to 0.
Therefore, determining how to use this method without damaging the instrument and how to control for the excess Cl 
will allow the study of the Cl/Br ratio with \ICPMS as well.

With this expansion of the trace element ratios we could study, several natural next steps present themselves.
Repeating the \XRF measurements to obtain Cl, Br, B, Ga, Sr, Ba, S would establish whether the
calibrated method yields results comparable to \ICPMS. 
The additional data for B/Ga would allow studying that additional ratio, also extensively discussed by \textcite{wei-2020}.
Repeating the \ICPMS measurements for Cl/Br would provide cross-validation for the XRF data.

That one previously published ratio requires local calibration, suggests that 
these ratios may not generalize as well as suggested by \textcite{wei-2020}. \Cref{fig:sr-ba-wei} shows that there
is significant variability in the dataset used by \textcite{wei-2020}. 
Whether this variability is caused by samples coming from different areas, or whether there can be such large
variability within well constrained areas is an important question regarding the generalizability of these results.

\subsection{Geology}

Between geology and geochemistry rests the challenge of independently classifying samples as
marine or freshwater sediments, e.g. with microfossils \parencite{kempf_lower_2005}.
This independent classification would enable assessing 
the precision and recall of the element ratios as classifiers, as well as creating a standard to calibrate the
ratio thresholds to local conditions.

With a validated calibration dataset and the corresponding thresholds, a study of an
extended area of the subalpine molasse would be feasible. Studying a larger area could provide
insights how frequent misclassified beds occur in the \USM and \UMM.
In any large study, a few of the samples taken should still be classified independently to ensure
the local calibration remains adequate.

The results of our study provide additional information about the stratigraphy
of the outcrops we sampled. Open questions remain regarding the origin of the marine layer JWL~1.
Fieldwork should establish the nature of the currently overgrown, underlying layer, and 
whether there are signs
of fault movement indicating it layer may have been moved into place, or if it indeed
indicates a possibly short transgression to marine conditions.
The latter would also suggest that the sediment was deposited in shallow conditions, which
could explain the presence of glauconite in some areas \parencite{Linth_Text,nichols-sedi-2023}.

Several additional steps should be taken here:
\begin{enumerate}
    \item Sample the beds on both sides of the gap for additional classifications.
    \item Find bedrock in the gap to complete the log, and sample some of the beds.
    \item Identify and describe fault movement, if indicators exist.
    \item Study other outcrops that belong to this stratigraphic sequence for more data.
\end{enumerate}